% Canonical prospective manuscript for workstream B.
% IMPORTANT: all numerical results are projected success values, not measured
% evidence. Replace them from frozen evaluation artifacts before circulation.
\documentclass[letterpaper,10pt,journal,twoside]{IEEEtran}

\usepackage{amsmath}
\usepackage{amssymb}
\usepackage{booktabs}
\usepackage{graphicx}
\usepackage{xcolor}
\usepackage{hyperref}
\usepackage{cite}

\hypersetup{colorlinks=true,linkcolor=blue,urlcolor=blue,citecolor=blue}
\newcommand{\method}{Wrench2Role}

\title{From Wrenches to Roles: Interaction-Grounded Object Segmentation\
in Open-World Robot Demonstrations}
\author{Anonymous Authors}

\begin{document}
\maketitle

\begin{center}
\fcolorbox{red}{red!5}{\parbox{0.94\columnwidth}{\small
\textbf{Prospective draft.} Numerical results are expected success values used
to design the completed paper. They are not measured results and must be
replaced after the frozen experiments.}}
\end{center}

\begin{abstract}
Robot demonstrations contain evidence not only of \emph{when} interaction
occurs, but of \emph{which} visible objects participate in it. Most visual
grounding pipelines nevertheless discard the gripper and force/torque signals
recorded alongside video. We introduce \method, a class-agnostic method that
assigns visual proposals the physical roles \emph{grasped object} and
\emph{contact receiver}. Gripper transitions and attachment motion ground the
grasped role, while a wrist wrench defines a 3-D line of action that is compared
with each mask's depth support. An uncertainty-aware ranker fuses interaction
and geometric evidence and abstains when the role is ambiguous; selected masks
are propagated with a frozen video segmentation model. We train and stress-test
the ranker in simulation with exact contact, instance, and transform ground
truth, then freeze it for evaluation on 36 real Franka demonstrations with
held-out objects, scenes, and recording sessions. In the projected successful
study, \method{} reaches 0.94/0.82 precision/coverage for grasped objects and
0.93/0.80 for contact receivers. It improves held-out simulated top-1 role
accuracy by 13.5 percentage points over the stronger single-modality baseline
and produces tracks with 0.84 mean IoU. Controlled perturbations identify when
wrench geometry remains informative and when the system should abstain. These
results establish measured interaction as a practical grounding signal for
open-world object-role segmentation in robot learning from demonstration.
\end{abstract}

\section{Introduction}

Learning from demonstration (LfD) seeks reusable task structure from
human-guided or kinesthetic robot executions. Constraint and interaction models
can infer \emph{how} an end effector moves or \emph{when} contact changes
\cite{overbeek2025frontiers,overbeek2025adrr}, but reproduction also requires a
visual answer to a more elementary question: which object was held, and which
object received the interaction? Semantic detectors are an awkward fit because
these are physical roles, not categories. A receiver may be a plate, tool,
fixture, articulated part, or previously unseen object.

Class-agnostic segmentation models can delineate almost any visible object, but
they do not decide which proposal matters. Motion-based discovery can associate
an unknown object with a gripper \cite{boerdijk2021whats}; human-interaction
methods similarly associate hands with held objects
\cite{higgins2023moves,narasimhaswamy2024hoist}. Contact-receiver identity is
harder. Several objects may be close to the tool, the relevant surface may be
occluded, and an eye-in-hand camera makes background motion an unreliable cue.
Conversely, force/torque sensing detects physical interaction but does not
directly identify an image region.

Our hypothesis is that these modalities resolve complementary ambiguities. A
gripper transition limits the grasped-object search temporally, and object
motion relative to the gripper supplies attachment evidence. During external
contact, the measured wrench defines a line of possible contact locations
\cite{bicchi1990}. Once transformed into the camera frame, this line can be
compared with the 3-D support of class-agnostic masks. Visual evidence resolves
the wrench's along-line ambiguity; force geometry resolves ambiguity among
visually plausible objects.

We instantiate this idea in \method{} (Fig.~\ref{fig:overview}). The system
detects interaction events, generates a common proposal set for every ablation,
computes role-specific interaction and wrench-geometry evidence, and selects or
rejects a candidate using calibrated confidence. Frozen SAM~2
\cite{ravi2024sam2} then propagates the mask bidirectionally. Simulation
provides exact labels and controlled corruptions; all thresholds are frozen
before real-world testing.

Our contributions are:
\begin{itemize}
  \item We formulate \emph{open-world object-role grounding}: assigning
  class-agnostic masks the roles of grasped object and contact receiver from
  synchronized eye-in-hand RGB-D, gripper state, pose, and wrist wrench.
  \item We introduce uncertainty-aware fusion of attachment evidence and the
  wrench line of action, with explicit abstention under ambiguous perception or
  geometry.
  \item We develop a simulation protocol with exact role ground truth and
  controlled perturbations of calibration, synchronization, sensing, clutter,
  and contact structure.
  \item We evaluate a frozen model on object-, scene-, and session-held-out real
  demonstrations against vision-only, interaction-only, and geometry-only
  alternatives.
\end{itemize}

\begin{figure}[t]
\centering
\includegraphics[width=\columnwidth]{figures/pipeline.png}
\caption{System overview. Events identify candidate moments, and all methods
share one class-agnostic mask pool. Role-specific visual, attachment, and
wrench-geometry evidence is fused with calibrated uncertainty. The system
selects a role mask for propagation or safely abstains. This placeholder will
be replaced by the final \method{} diagram.}
\label{fig:overview}
\end{figure}

\section{Related Work}

\subsection{Interaction-grounded object segmentation}

``What's This?'' learns unknown object-in-gripper segmentation from robot
motion without manual object labels \cite{boerdijk2021whats}. MOVES learns
pixel grouping and hand--object association from manipulation video
\cite{higgins2023moves}; HOIST-Former jointly identifies, segments, and tracks
hand-held objects \cite{narasimhaswamy2024hoist}. HOI-DETR extends human-object
interaction detection to a hand, first object, and second acted-upon object
\cite{darkhalil2026hoi}. These works establish interaction as a powerful visual
prior. Our camera instead moves with the robot, roles are inferred from robot
measurements, and external contact must be assigned to an object instance.

SAM \cite{kirillov2023sam} and SAM~2 \cite{ravi2024sam2} offer strong
class-agnostic masks given a prompt. We use them as proposal and tracking
infrastructure. Our problem is choosing the physically relevant proposal
without an object name or recording-specific manual prompt.

\subsection{Visual and force-based contact localisation}

Im2Contact predicts dense extrinsic contact from a calibrated RGB-D view and
robot proprioception without tactile or force sensing \cite{kim2023im2contact}.
It demonstrates that task-aligned visual and end-effector cues can localise
diverse contact. We instead ask which open-world instance occupies each
physical role and explicitly incorporate a measured wrist wrench. The methods
are complementary; neither motivates a claim that vision or force is universally
sufficient.

Intrinsic contact sensing constrains a single contact from a wrist wrench
\cite{bicchi1990}. Bimbo \emph{et al.} use force/torque to infer possible
contacts and update 3-D occupancy maps \cite{bimbo2022mapping}. Our output is
instead an image-space, temporally propagated instance mask with a physical
role. Hoi! provides large-scale synchronized visual, force, and tactile data
across embodiments \cite{engelbracht2026hoi}; it motivates external evaluation
but benchmarks a distinct articulated-interaction setting.

\section{Problem Formulation}

A demonstration is a synchronized sequence
\begin{equation}
\mathcal{D}=\{I_t,D_t,\mathbf{T}^{W}_{E,t},g_t,
\mathbf{w}_t\}_{t=1}^{T},
\end{equation}
where $I_t,D_t$ are eye-in-hand RGB and depth,
$\mathbf{T}^{W}_{E,t}$ is end-effector pose, $g_t$ is gripper state, and
$\mathbf{w}_t=[\mathbf{f}_t^\top,\boldsymbol{\tau}_t^\top]^\top$ is wrist
wrench. At event $e$, a class-agnostic segmenter produces candidates
$\mathcal{M}_e=\{m_i\}_{i=1}^{N_e}$. We predict
$m_e^{r}\in\mathcal{M}_e\cup\{\varnothing\}$ for
$r\in\{\mathrm{grasped},\mathrm{receiver}\}$; $\varnothing$ denotes
abstention. Category labels and object templates are unavailable.

Evaluation separates candidate recall, role selection, and tracking. Proposal
recall asks whether any $m_i$ overlaps the reference object. Selection measures
top-$k$ role accuracy and precision as a function of automatic coverage.
Tracking is evaluated only after successful selection.

\section{Method}

\subsection{Event-conditioned proposals}

Grasp and release events are detected from gripper-width transitions after a
minimum-travel test rejects noise. Contact candidates are peaks in gravity- and
bias-compensated wrench magnitude, grouped into activity clusters and restricted
to the active grasp interval when one exists. A short temporal window is
retained rather than committing to one frame.

At each anchor, a frozen class-agnostic segmenter produces masks at multiple
granularities. Near-duplicates are suppressed by mask IoU before confidence
margins are computed. Depth validity, border truncation, and minimum support
tests remove uninterpretable candidates; no category or recording-specific
pixel prior is used.

\subsection{Grasped-object evidence}

For candidate $m_i$, the grasp descriptor $\boldsymbol{\phi}^{G}_i$ combines
overlap with the projected gripper; change in support across closure; motion
coherence with the end effector; residual motion relative to the background;
and persistence during early hold. Persistence alone cannot promote a mask: a
static fixture may be stable yet unrelated to the grasp. Features are aggregated
across the event window using robust medians and dispersion estimates.

\subsection{Wrench-grounded contact evidence}

Under a single-resultant-contact model, force $\mathbf{f}$ and torque
$\boldsymbol{\tau}$ define the closest point of the line of action to the sensor
origin and its direction \cite{bicchi1990}:
\begin{equation}
\mathbf{r}_0=\frac{\mathbf{f}\times\boldsymbol{\tau}}
{\lVert\mathbf{f}\rVert^2},\qquad
\hat{\mathbf{f}}=\frac{\mathbf{f}}{\lVert\mathbf{f}\rVert}.
\label{eq:line}
\end{equation}
The hand--eye transform and time-varying pose map this line into the camera
frame. Valid depth pixels within $m_i$ yield $\mathcal{P}_i$. Geometric
discrepancy is
\begin{equation}
d_i=\min_{\mathbf{p}\in\mathcal{P}_i,\,s\in\mathbb{R}}
\|\mathbf{p}-(\mathbf{r}_0+s\hat{\mathbf{f}})\|_2.
\label{eq:distance}
\end{equation}
The receiver descriptor $\boldsymbol{\phi}^{C}_i$ includes $d_i$, mask support
near the line, depth continuity, force direction relative to the local surface,
and temporal consistency around contact.

A single projected line is overconfident under calibration, synchronization,
and wrench error. We sample their estimated joint uncertainty, recompute
Eq.~\eqref{eq:distance}, and retain its expectation and dispersion. Masks score
highly only when their support remains plausible across likely sensing error.

\subsection{Fusion, abstention, and propagation}

A shallow role-specific network fuses visual/interaction evidence, geometry,
and missing-modality indicators:
\begin{equation}
z_i^r=h_r(\boldsymbol{\phi}_i^G,\boldsymbol{\phi}_i^C,
\boldsymbol{\phi}_i^V),\qquad
p_i^r=\operatorname{softmax}_{i}(z_i^r).
\end{equation}
Training combines cross-entropy with pairwise ranking against hard negatives
from the same frame. Temperature scaling on disjoint validation data calibrates
$p_i^r$. The system accepts the top mask only if its probability and margin
over the best non-duplicate candidate exceed frozen thresholds. Otherwise it
abstains.

Accepted masks seed frozen SAM~2 bidirectional propagation. Role, confidence,
provenance, and per-frame masks are written to a versioned sidecar for
downstream LfD.

\section{Experimental Protocol}

\subsection{Simulation and stress tests}

Isaac Sim reproduces the Franka arm, gripper, eye-in-hand camera, and wrist
sensor geometry. Scripted grasp--carry--contact episodes randomize object shape,
distractors, lighting, texture, mass, friction, approach direction, and contact
type. The simulator exports synchronized RGB, depth, pose, gripper, wrench,
instance masks, contact pairs, contact points, and exact role labels.

The prospective study contains 1,200/200/400 train/validation/test episodes.
Test objects and scene assets never occur in training. Three independent
generation and training seeds are used. Separate suites perturb extrinsic
rotation and translation, time offset, wrench bias/noise, depth dropout,
occlusion, clutter, and multiple contacts. Before generation, the virtual wrench
is checked in ten deterministic scenes against analytically known loads.

\subsection{Real-robot study}

The real platform comprises a Franka Research~3, Franka Hand, eye-in-hand ZED
Mini RGB-D camera, and Bota SensONE wrist force/torque sensor. A ChArUco hand--eye
procedure estimates the camera-to-measurement transform. Frame names, direction,
units, reprojection residuals, and independent physical checks are recorded
before evaluation.

After a separate pilot, the projected main dataset contains 36 independent
demonstrations across 12 held/receiver object pairings, six scene layouts, and
four recording sessions. Entire object identities and sessions are held out.
Reference masks and roles are authored without method scores. Model, thresholds,
and abstention policy are frozen after simulation and validation data.

\subsection{Baselines and metrics}

All internal baselines receive identical proposals: a region/size heuristic;
DINOv2 \cite{oquab2023dinov2} attention with depth; a vision/interaction-only
ranker; wrench geometry alone; and the full hybrid. DistinctNet and HOI-DETR are
reported as out-of-domain transfer baselines where their outputs map fairly to
robot roles. Receiver evaluation also includes an Im2Contact-style
visual/end-effector baseline.

Primary metrics are top-1 accuracy, precision--coverage curves, coverage at
0.90 precision, and selective risk. We also report proposal recall, seed IoU,
track IoU, 3-D contact discrepancy, and corruption sensitivity. Confidence
intervals are clustered by session, never by video frame.

\section{Expected Results}

\subsection{Held-out role selection}

Table~\ref{tab:real} shows the projected successful result. Attachment evidence
substantially improves grasped-object ranking. Geometry alone is precise for
contact receivers but abstains when depth or calibration uncertainty leaves
several candidates near the line. Fusion is the only method expected to exceed
the preregistered 0.90 precision and 0.75 coverage bar for both roles.

\begin{table}[t]
\centering
\caption{\textbf{Projected, not measured} held-out real selection. P/C is
precision/coverage at the frozen operating point.}
\label{tab:real}
\setlength{\tabcolsep}{3.5pt}
\begin{tabular}{lcccc}
\toprule
& \multicolumn{2}{c}{\textbf{Grasped}} &
\multicolumn{2}{c}{\textbf{Receiver}} \\
\cmidrule(lr){2-3}\cmidrule(lr){4-5}
\textbf{Method} & \textbf{Top-1} & \textbf{P/C} &
\textbf{Top-1} & \textbf{P/C} \\
\midrule
Region/size & .44 & .61/.36 & .39 & .57/.31 \\
DINOv2 + depth & .53 & .68/.43 & .48 & .64/.39 \\
Vision/interaction & .81 & .88/.72 & .70 & .84/.63 \\
Wrench geometry & -- & -- & .76 & .91/.68 \\
\method{} & \textbf{.89} & \textbf{.94/.82} &
\textbf{.87} & \textbf{.93/.80} \\
\bottomrule
\end{tabular}
\end{table}

On object-and-scene-held-out simulation, the hybrid reaches a projected 89.3\%
top-1 role accuracy, compared with 72.4\% for vision/interaction and 75.8\% for
geometry. Its 13.5-point gain over the stronger component exceeds the
predetermined continuation gate. Correct role masks occur in the proposal pool
in a projected 96.8\% of events, leaving proposal failure as a measurable but
secondary source of error.

\subsection{Tracking and geometric accuracy}

Accepted seeds achieve a projected mean IoU of 0.88; bidirectional propagation
produces 0.84 mean track IoU while the object is visible. Explicit absence
handling prevents persistent mask growth after an object leaves view. The
projected median 3-D distance from the estimated line of action to the reference
contact region is 1.7\,cm, with a 90th percentile of 3.8\,cm.

\subsection{Ablation and robustness}

Removing wrench geometry causes the largest receiver degradation; removing
attachment evidence primarily affects grasped objects. Removing uncertainty
sampling changes clean-condition accuracy little but makes confidence poorly
calibrated under transform and timing errors. At a projected $5^\circ$
extrinsic rotation error, uncertainty-aware fusion retains 0.79 coverage at
0.90 precision, compared with 0.58 for deterministic fusion. At $10^\circ$,
both degrade, but the uncertainty-aware system abstains rather than maintaining
a falsely high acceptance rate. Multiple simultaneous contacts remain the
dominant physical failure because a single resultant wrench need not identify
either contact independently.

\section{Discussion}

The expected result does not imply that vision alone is incapable of contact
understanding. Im2Contact and modern HOI systems demonstrate the opposite. The
narrow conclusion is that measured interaction supplies complementary evidence
for assigning open-world object \emph{roles}, especially when several masks are
visually plausible. A common proposal pool and modality ablations test this
without conflating proposal quality with role selection.

Simulation supplies exact supervision and controlled perturbations, but cannot
reproduce every Bota artifact, ZED failure, compliance effect, or latency.
Simulation therefore establishes mechanism and sensitivity; held-out physical
recordings establish transfer. Adjacent frames and duplicate exports are never
treated as independent trials.

The method assumes a dominant resultant contact and sufficient valid depth near
a proposal. Transparent, reflective, and occluded objects can violate the
latter; distributed or simultaneous contacts violate the former. Calibration
uncertainty can be modelled but not made irrelevant. Such cases motivate
abstention and are reported as reduced coverage rather than excluded failures.

\section{Conclusion}

We presented \method, a selective grounding method that converts robot
interaction measurements into prompts for open-world video segmentation.
Gripper transitions and attachment motion identify the grasped role; calibrated
wrench geometry constrains the receiver; uncertainty-aware fusion resolves
ambiguities remaining in either modality. If reproduced by the frozen
experiments, the projected gains support a practical principle for robot LfD:
the robot's physical interaction is not merely a control signal, but a grounding
signal for visual object identity.

\begin{thebibliography}{99}
\bibitem{bicchi1990} A.~Bicchi, ``Intrinsic contact sensing for soft fingers,''
in \emph{Proc. IEEE Int. Conf. Robot. Autom. (ICRA)}, 1990.
\bibitem{kirillov2023sam} A.~Kirillov \emph{et al.}, ``Segment Anything,'' in
\emph{Proc. IEEE/CVF Int. Conf. Comput. Vis. (ICCV)}, 2023.
\bibitem{ravi2024sam2} N.~Ravi \emph{et al.}, ``SAM 2: Segment Anything in
images and videos,'' in \emph{Proc. Int. Conf. Learn. Represent.}, 2025.
\bibitem{oquab2023dinov2} M.~Oquab \emph{et al.}, ``DINOv2: Learning robust
visual features without supervision,'' \emph{TMLR}, 2024.
\bibitem{boerdijk2021whats} W.~Boerdijk, M.~Sundermeyer, M.~Durner, and
R.~Triebel, ``What's This?---Learning to segment unknown objects from
manipulation sequences,'' in \emph{Proc. ICRA}, 2021.
\bibitem{higgins2023moves} R.~E.~L. Higgins and D.~F. Fouhey, ``MOVES:
Manipulated objects in video enable segmentation,'' in \emph{Proc. CVPR}, 2023.
\bibitem{narasimhaswamy2024hoist} S.~Narasimhaswamy, H.~A. Nguyen, L.~Huang,
and M.~Hoai, ``HOIST-Former: Hand-held objects identification, segmentation,
and tracking in the wild,'' in \emph{Proc. CVPR}, 2024.
\bibitem{darkhalil2026hoi} A.~Darkhalil, D.~Damen, and D.~Fouhey, ``Improving
and evaluating hand-object interaction detection,'' \emph{arXiv:2606.17384},
2026.
\bibitem{kim2023im2contact} L.~Kim, Y.~Li, M.~Posa, and D.~Jayaraman,
``Im2Contact: Vision-based contact localization without touch or force
sensing,'' in \emph{Proc. CoRL}, 2023, pp.~1533--1546.
\bibitem{bimbo2022mapping} J.~Bimbo, A.~S. Morgan, and A.~M. Dollar,
``Force-based simultaneous mapping and object reconstruction for robotic
manipulation,'' \emph{IEEE Robot. Autom. Lett.}, vol.~7, no.~2,
pp.~4749--4756, 2022.
\bibitem{engelbracht2026hoi} T.~Engelbracht \emph{et al.}, ``Hoi!---A
multimodal dataset for force-grounded, cross-view articulated manipulation,''
in \emph{Proc. CVPR}, 2026, pp.~8880--8890.
\bibitem{overbeek2025frontiers} A.~H.~G. Overbeek, D.~Dresscher,
H.~van~der~Kooij, and M.~Vlutters, ``Versatile kinematics-based constraint
identification applied to robot task reproduction,'' \emph{Frontiers Robot.
AI}, vol.~12, art.~1574110, 2025.
\bibitem{overbeek2025adrr} A.~H.~G. Overbeek, H.~van~der~Kooij, D.~Dresscher,
and M.~Vlutters, ``Identifying physical interactions in contact-based robot
manipulation for learning from demonstration,'' \emph{Advanced Robotics
Research}, e202500109, 2025.
\end{thebibliography}

\end{document}
